\label{sec:verification}

\subsection{The Verification Protocol}

The SHA-256 verification protocol provides cryptographic proof that
\textsc{bisect-apportion}'s Huntington-Hill output matches the
Census Bureau's published 2020 apportionment.
The protocol:

\begin{enumerate}
  \item \textbf{Construct the canonical string}: Sort the 50 (state, seats)
    pairs alphabetically by state name.
    Concatenate as \texttt{"StateName:seats\textbackslash n"}
    for each pair (e.g., \texttt{"Alabama:7\textbackslash nAlaska:1\textbackslash n...Wyoming:1\textbackslash n"}).
  \item \textbf{Compute SHA-256}: Apply SHA-256 to the canonical string
    encoded as UTF-8.
  \item \textbf{Compare to reference hash}: The Census Bureau table
    \citep{census2020apportionment} is processed into the same canonical
    format, and its SHA-256 is committed to the repository as a
    compile-time constant.
  \item \textbf{Verify}: If the computed and reference hashes match,
    the apportionment is correct; otherwise, the verification fails
    and the CI pipeline exits non-zero.
\end{enumerate}

\begin{verbatim}
pub fn verify_sha256(
    allocation: &[(String, usize)],
    reference_hash: &str,
) -> bool {
    let mut canonical = String::new();
    let mut sorted = allocation.to_vec();
    sorted.sort_by_key(|(name, _)| name.clone());
    for (name, seats) in &sorted {
        canonical.push_str(&format!("{}:{}\n", name, seats));
    }
    let computed = sha256(canonical.as_bytes());
    computed == reference_hash
}
\end{verbatim}

\subsection{The 2020 Verification Result}

Running:
\begin{verbatim}
bisect apportion --year 2020 --method huntington-hill --verify
\end{verbatim}
produces:
\begin{verbatim}
Apportionment: 2020 Census, Huntington-Hill, 435 seats
SHA-256: [committed hash]
Verification: PASS (matches Census Bureau 2020 apportionment)
States verified: 50/50
\end{verbatim}

The committed reference hash corresponds to the Census Bureau's
April 26, 2021 apportionment table \citep{census2020apportionment}.

\subsection{2010 and 2000 Verification}

The same verification is performed for the 2010 and 2000 Census years.
Both pass with zero discrepancies across all 50 states.
This three-decade verification establishes that the implementation
is correct not only for 2020 but is robust to population distribution
changes across census years.

\subsection{Why SHA Verification Matters for Redistricting}

In redistricting litigation and commission proceedings, the
apportionment step is often taken for granted.
Parties arguing over district maps rarely question the seat counts
assigned to each state.
But if the redistricting platform produces incorrect seat counts,
every downstream computation is wrong.

SHA verification provides:
\begin{itemize}
  \item \textbf{Reproducibility}: Any party can reproduce the
    verification by running \texttt{bisect apportion --verify}.
  \item \textbf{Audit trail}: The hash comparison creates a documented
    link between the Census Bureau's official publication and the
    redistricting platform's input.
  \item \textbf{CI enforcement}: The verification runs in CI on every
    commit, ensuring that code changes do not silently break the
    apportionment computation.
\end{itemize}
